% III. Проектиране на системата.
% Слоеве, отговорности, граници между тях. Фигурите се добавят, когато схемите
% бъдат начертани; дотогава мястото им е маркирано с \TODO.
\section{Проектиране на системата}
\label{sec:design}

\subsection{Слоеве и отговорности}

Системата е разделена на пет слоя с еднопосочна зависимост, отдолу нагоре. Най-отдолу
стои \term{транспортът}, който притежава сокета и буферите и не знае нищо за
съдържанието на байтовете. Над него стои \term{протоколният кодек}, по един за
PostgreSQL и за MySQL, който превръща поток от байтове в съобщения и обратно. Следва
\term{декодерът на редове}, който превръща тялото на едно съобщение с данни в изгледи
към стойности, без да ги копира. Над него е \term{пулът от връзки}, който управлява
живота на връзките и кеша от подготвени заявления. Най-отгоре стои \term{слоят за
типово изображение}, който свързва стойностите от протокола с типовете на C++.

Границата между всеки два слоя е тесен договор, а не общ обект със състояние. Целта е
всеки слой да се тества самостоятелно: кодекът да се захранва със записан байтов
поток, декодерът с готово тяло на съобщение, пулът с фиктивен транспорт. Това
разделение е и условието за смисленото сравнение с доставчиковите библиотеки, защото
позволява цената да се отнесе към конкретен слой.

\subsection{Протоколни кодеци}

Кодекът на PostgreSQL следва съобщителния модел на протокола: етикет, дължина, тяло, и
двата цикъла на изпълнение, простия и разширения~\cite{pg_protocol}. Кодекът на MySQL
следва пакетния модел с номер на последователност и цели числа с променлива
дължина~\cite{mysql_protocol}. Двата кодека нямат общ базов клас. Общото между тях е
изразено като изисквания към типа, тоест като \code{concept}, което позволява
останалите слоеве да са шаблонни и да не плащат за динамично изпращане.

\TODO{фигура: диаграма на слоевете и на потока от данни между сокета и типизираната
стойност; да се начертае с TikZ и да се озаглави по правилото за фигури}

\subsection{Декодиране без копие}

Резултатният ред пристига като част от получаващия буфер, притежаван от транспорта.
Декодерът не изгражда нови обекти за стойностите, а издава изгледи от вида
\code{std::span<const std::byte>} и \code{std::string\_view} към този буфер. Копие се
прави само когато приложението изрично поиска притежание на стойността. Условието за
коректност е ясно: изгледите са валидни, докато редът е жив, и това
ограничение е част от публичния договор, а не подразбираща се подробност.

Тази схема прехвърля отговорност върху приложението и затова е рискова. Проектът я
приема, но я подсигурява с проверки при активирани диагностики и с тестове, които
целенасочено използват изглед към вече невалиден ред, за да се убедят, че диагностиката го
улавя. \TODO{да се опише точният механизъм за проверка след като бъде избран}

\subsection{Пул от връзки и кеш на подготвените заявления}

Пулът е асинхронен. Заявка за връзка, която не може да бъде удовлетворена веднага,
спира изпълнението на извикващата съпрограма, вместо да блокира нишка. Всяка връзка
носи собствен кеш на подготвените заявления, защото подготвеното заявление в двата
разглеждани протокола е свързано със сесията, а не с базата. Кешът е с ограничен
размер и с политика на изместване. \TODO{избор на политика и на подразбиращ се размер,
след като измерванията покажат чувствителността към тях}

\subsection{Слой за типово изображение}

Слоят превръща стойност от протокола в типизирана стойност на C++ и обратно, като
съответствието се задава декларативно и се проверява при компилация. Целта е
несъответствие между схемата и заявената в кода структура да бъде грешка при
компилация или ясна грешка при изпълнение, а не мълчаливо погрешно четене на байтове.
\TODO{таблица със съответствието между типовете на PostgreSQL, типовете на MySQL и
типовете на C++, включително случаите без точно съответствие}
